\originalchapter{淬火无序、聚合物与大纲补充}
\label{ch:random}
\source{8.334大纲列随机体系与聚合物；26份公开讲义未为这些大纲项各给独立全讲，本章为编者补写的基础连接，不冒称遗漏的课堂原稿。}
\originalsection{退火与淬火的区别}
若杂质随系统充分运动，可与热变量一起求和，退火自由能$F_a=-\kb T\log\overline Z$；若杂质在观测时间内固定，每个样本先取热自由能，淬火$F_q=-\kb T\,\overline{\log Z}$。由$\log$凹性$\overline{\log Z}\le\log\overline Z$，故$F_q\ge F_a$。平均哈密顿量、平均配分函数、平均自由能不是同一运算，不能互换。

例如独立随机能级$E$的单态$Z=\e^{-\beta E}$，淬火$F_q=\overline E$，退火$F_a=\overline E-\beta\operatorname{Var}E/2$对Gaussian成立。退火结果被罕见低能样本偏重，差异不是仅一个数值误差。
\originalsection{Harris与Imry--Ma判据}
短程独立随机局域临界温度在相关体积$\xi^d$内的平均涨落$\delta t_\xi\sim\xi^{-d/2}$。洁净体系距离临界$t\sim\xi^{-1/\nu}$。若随机涨落比$t$更小，要求$d\nu>2$；反之$d\nu<2$无序相关。边际等号需更细RG，不能据此判定必然相关或无关。若洁净超尺度成立，等价$\alpha<0$使弱随机质量无关。

随机场耦合$-\int h(x)\phi(x)$则不同。Ising翻一个尺度$L$畴的壁代价$\sim\sigma L^{d-1}$，随机场能增益$\sim hL^{d/2}$。$d<2$时大畴增益压过壁，低维长程序被破坏；$d=2$边际需更精细分析。连续自旋可通过缓慢转向降低代价为$L^{d-2}$，比较给下临界尺度4。随机键与随机场不能用同一个Harris不等式混称。
\originalsection{副本与自旋玻璃的范围}
副本恒等式形式$\overline{\log Z}=\lim_{n\to0}(\overline{Z^n}-1)/n$，对正整数$n$可先计算$n$份热变量共享同一无序。若随机场Gaussian，平均$\exp[\beta\int h\sum_a\phi_a]$生成$\frac12\beta^2\Delta\int(\sum_a\phi_a)^2$，使不同副本耦合。随机质量则生成四次副本耦合。这解释副本作用量项的来源。

从整数$n$解析继续到0并非无需证明的普遍合法步骤；副本对称鞍点也可能失稳。自旋玻璃以无序冻结局域磁矩与大量亚稳态为特征，可用Edwards--Anderson量$q_{EA}=\overline{\avg{s_i}^2}$区别均匀磁化为0的冻结态。完整有限维自旋玻璃与副本破缺的严谨解不属于这组公开讲义的可重建证明，本章给定义、平均结构与判据，不以一个$q$方程宣称完成所有模型。
\originalsection{随机链与Flory尺度}
理想柔性链有$N$段、段长$a$，独立方向使$\avg{R^2}=Na^2$，Gaussian端点分布$P(R)\propto\exp[-dR^2/(2Na^2)]$。排斥体积$ v>0$可近似自由能$F/\kb T\sim R^2/(Na^2)+vN^2/R^d$，前项是伸展熵代价，后项为单体密度平方积分。极小$R^{d+2}\sim v a^2N^3$，Flory指数$\nu_F=3/(d+2)$，三维为$3/5$。

这是量级平衡而非精确指数；它忽略链段相关，二维虽恰得$3/4$也不能据此称该方法严格。$d>4$排斥体积无关，应回理想$\nu=1/2$，而不是无条件继续$3/(d+2)$。$n\to0$的$O(n)$场论把两点高温图中一条开链保留、闭环因$n$权重消失，建立自避行走与RG的联系；近四维前章$\nu=1/2+\epsilon/16+O(\epsilon^2)$给可控修正。

吸引与排斥抵消的$\theta$点涉及三体稳定项与三临界，上临界维数3。更复杂网络需要按节点功能与连接拓扑另立熵计数，不能把所有网络都视为单条线性链。课程大纲的这些扩展保留实际覆盖与证明边界。
